\section{アプローチ06（最小入力サイズの拡大・保証手法）}

\subsection{手法の概要}
アプローチ05で発生した「一次検出で得られたバウンディングボックスの領域が小さすぎる場合に、ポーズ推定（MoveNet Lightning）の入力解像度を満たせず精度が低下する」という課題に対応するため、アプローチ06では入力画像切り出し直後のポーズ推定前段階において、切り取られた画像がモデルの最小推奨入力サイズを下回る場合に自動拡大・サイズ保証を行う前処理を導入した。

このサイズ拡大処理は、メインスレッド（\texttt{detectPeople.ts} 内の \texttt{Promise.all} 子処理）ではなく各 Worker 内で実施し、固定キャンバスの保持数を最適化した。具体的な設計および最適化仕様は以下の通りである。

\begin{itemize}
    \item \textbf{Worker内固定キャンバスの流用}：Workerごとに、小さすぎるビットマップ画像を受信した際のリサイズ用固定キャンバスをグローバルスコープで保持・使い回し、頻繁な DOM キャンバス生成によるメモリ消費を抑制した。
    \item \textbf{拡大基準と十分な解像度確保}：固定キャンバスのサイズは元画像の解像度ではなく、モデル処理において十分な精度が担保できる定格サイズを流用した。切り出されたビットマップがすでに十分なサイズを有している場合は拡大を行わず、そのまま推定に回す設計とした。
    \item \textbf{座標系の復元・拡大比率計算}：拡大された画像領域上で取得されたポーズ検出結果（各関節の二次元座標 $P_{scaled} = (x_{s}, y_{s})$）に対し、元のバウンディングボックスサイズ $W_{orig} \times H_{orig}$ と拡大後サイズ $W_{target} \times H_{target}$ から求めたスケール比率を用いて、元の画像空間上の座標 $P_{orig}$ へと逆算・正規化を行った。
\end{itemize}

\subsection{結果およびアプローチ05との比較・課題}

\subsubsection{アプローチ05との比較}
小さく映る遠方の人物や解像度の低いバウンディングボックス領域であっても、モデル推奨サイズまで引き延ばして推論を行うことで、MoveNet Lightning における関節キーポイントの検出成功率および信頼度（confidence score）が著しく向上した。これにより、待機列後方の小規模な人物領域に対するポーズ検出漏れが解消された。

\subsubsection{判明した課題}
\begin{enumerate}
    \item \textbf{重複検出バウンディングボックスに対する余剰計算}：
    低信頼度で列挙された多数の候補それぞれに対してリサイズおよびポーズ推定を実行するため、同じ人物の重なり合った複数領域に対して過剰な計算資源を消費する問題が依然として残った。
    \item \textbf{枠ずれに伴うノイズキーポイントの発生}：
    バウンディングボックスの切り出し位置が本来の人物中心からずれている場合、拡大によって背景ノイズが強調され、誤った姿勢ベクトルが算出されるリスクが生じた。
\end{enumerate}
